\clearpage\begingroup
\HistoricalHead{186}{L}{PHYSIK.\quad NACHLASS.}
\begin{center}8.\end{center}
\noindent\hspace*{\parindent}Der Multiplicator, welcher die grosse Nadel im magnetischen Observatorium umgibt, besteht aus zwei ganz gleichen Hälften. Jede Hälfte enthält 135 Umwindungen, die in drei durch Seidenband von einander getrennten Schichten einen rechtwinklichten hölzernen Rahmen umgeben. Jede Schicht enthält daher 45 Umwindungen, die schraubenförmig zwischen zwei Grenzrechtecken laufen, die 67 Millimeter von einander abstehen. Es ist daher $k'-k^0=67^{\mathrm{mm}}$ und $\lambda=1^{\mathrm{mm}}{,}4889$. Es ist zureichend die Rechnung für die mittelste Schicht zu führen, für welche die Länge des Rechtecks $=1341^{\mathrm{mm}}$, die Breite oder vielmehr Höhe, da die Rechtecke vertical stehen, $=206^{\mathrm{mm}}$. Die beiden Rahmen stehen symmetrisch gegen den magnetischen Meridian auf der Nordseite dicht zusammen, während sie auf der Südseite etwas von einander abstehen, um die Spiegel am Südpol der Nadel frei zu lassen; jeder Rahmen macht mit dem magnetischen Meridian einen Winkel von $2^0\,9'\,19''=\theta$. Die Länge der Nadel beträgt $1222^{\mathrm{mm}}{,}4$; der[*)]\begingroup\renewcommand{\thefootnote}{}\footnotetext{[*) In der Handschrift steht nach mehreren Durchstreichungen statt \textls[120]{der} das Wort »ihr«.]}\endgroup\ Mittelpunkt und die Axen des östlichen und westlichen Rahmens treffen in dem Mittelpunkte der Nadel zusammen; der Abstand dieses Mittelpunktes von dem nächsten Rechtecke auf jeder Seite beträgt $32^{\mathrm{mm}}{,}5$.

Für diesen Mittelpunkt ist daher in Beziehung auf die mittlere Schicht jeder Multiplicatorhälfte in Millimetern
\[
\alpha=\gamma=670{,}5;\quad\beta=\delta=103;\quad k^0=32{,}5;\quad k'=99{,}5.
\]
\clearpage\SourceLeaf{192}\HistoricalHead{187}{R}{WECHSELWIRKUNGEN ZWISCHEN GALVANISCHEN STRÖMEN UND MAGNETISMUS.}
\noindent Die ganze Kraft normal gegen den Multiplicator, d. i. gegen die Ebene der Grenzrechtecke, wird
\[
=\frac{4}{\lambda}\left(\operatorname{Arc\,tg}\frac{k'e'}{\alpha\beta}-\operatorname{Arc\,tg}\frac{k^0e^0}{\alpha\beta}\right),
\]
\noindent während die mit den Seiten der Rechtecke parallelen verschwinden. Jene Formel gibt $\frac{4}{\lambda}(44^0\,38'\,55''-17^0\,43'\,25'')$, oder für alle drei Schichten $\frac{12}{\lambda}(26^0\,55'\,30'')$ oder wenn man den Bogen in Theilen des Halbmessers ausdrückt und für $\lambda$ seinen Werth setzt, $3{,}8731$. Dieser Werth gilt für jede Multiplicatorhälfte; allein für die eine Hälfte ist die Kraft dem Rahmen zugekehrt, für die andere davon abgekehrt, und zwar werden, wenn der positive Strom oben von Süden nach Norden den Multiplicatordraht durchläuft, beide Richtungen (in Beziehung auf das positive magnetische Fluidum) von Ost nach West gehen, aber für die westliche Hälfte mit einer Abweichung von $2^0\,9'\,19''=\theta$ südlich, für die östliche Hälfte mit einer eben so grossen Abweichung nördlich (alles wie sich von selbst versteht in Beziehung auf den \textls[120]{magnetischen} Meridian). Die ganze Kraft ist daher rein westlich, und ihr Werth wird gefunden, wenn man die obige Zahl mit $2\cos\theta$ multiplicirt, $=7{,}5697$.

Für denjenigen Punkt der durch das Centrum der Nadel gehenden magnetischen Axe, welcher am nördlichen Ende der Nadel liegt, haben wir
\[
\alpha=59{,}75;\quad\gamma=1281{,}25;\quad\beta=\delta=103;\quad k^0=9{,}5;\quad k'=76{,}5,
\]
\noindent und es gehen aus der Wirkung der mittleren Schicht jeder Multiplicatorhälfte zwei Kräfte hervor, eine senkrecht gegen den Rahmen, die andere parallel mit dessen horizontalen Seiten. Die erste Kraft
\[
\begin{aligned}
&=\frac{2}{\lambda}\left(\operatorname{Arc\,tg}\frac{k'e'}{\alpha\beta}+\operatorname{Arc\,tg}\frac{k'g'}{\beta\gamma}-\operatorname{Arc\,tg}\frac{k^0e^0}{\alpha\beta}-\operatorname{Arc\,tg}\frac{k^0g^0}{\beta\gamma}\right)\\[5pt]
&=\frac{2}{\lambda}(60^0\,23'\,8''+36^0\,44'\,20''-10^0\,26'\,51''-5^0\,17'\,12'')=\frac{2}{\lambda}\mathbin{.}81^0\,23'\,25''.
\end{aligned}
\]
\noindent Von der Richtung dieser Kraft gilt dasselbe, was vorher von dem Mittelpunkte bemerkt wurde, so dass aus allen drei Schichten beider Multiplicatorhälften die rein westlich wirkende Kraft $=\frac{12\mathbin{.}81^0\,23'\,25''}{\lambda}\cos\theta$ hervorgeht, also nachdem der Bogen in Theilen des Halbmessers ausgedrückt ist $=11{,}4408$. Die zweite Kraft wird für die mittelste Schicht jeder Multiplicatorhälfte ihrer Grösse nach [durch]
\[
\frac{2}{\lambda}\log\frac{\beta+e'}{\beta+e^0}\mathbin{.}\frac{\beta+g^0}{\beta+g'}
\]
\HistoricalFoot{}{24\textsuperscript{*}}
\clearpage\SourceLeaf{193}\HistoricalHead{188}{L}{PHYSIK.\quad NACHLASS.}
\noindent ausgedrückt. Obige Zahlen geben die briggischen Logarithmen von
\[
\begin{array}{lr}
\beta+e^0&2{,}34724\\
\beta+e'&2{,}38834\\
\beta+g^0&3{,}14252\\
\beta+g'&3{,}14322.
\end{array}
\]
\noindent Also Werth der Formel $=\frac{2}{\lambda}\mathbin{.}0{,}04040\mathbin{.}2{,}3025$, oder für alle drei Schichten
\[
=\frac{6}{\lambda}\mathbin{.}0{,}04040\mathbin{.}2{,}3025=0{,}3748.
\]
\noindent Die Richtung dieser Kraft ist aber (bei der oben angenommenen Richtung des Stromes) für den westlichen Theil des Multiplicators von Nord nach Süd, für den östlichen von Süd nach Nord, für beide mit der westlichen Abweichung $=\theta$. Man sieht daraus leicht, dass wenn man jede dieser Kräfte in zwei andere zerlegt, die eine parallel mit dem magnetischen Meridian, die andere senkrecht dagegen, die ersteren sich aufheben, so dass die Gesamtwirkung rein nach West gerichtet ist; ihre Grösse wird gefunden, wenn man den vorigen Werth mit $2\sin\theta$ multiplicirt, und ergibt sich $=0{,}0282$.

Die ganze Wirkung des Multiplicators in dem nördlichsten Punkte der Centrallinie der Nadel besteht also in einer Kraft nach Westen $=11{,}4690$.

Wir haben eine ähnliche Rechnung noch für 19 andere Punkte in derselben Linie, bis zum südlichen Endpunkte hin, geführt. Hier ist die Übersicht der Resultate, wobei mit 0 der nördliche, mit 20 der südliche Endpunkt bezeichnet ist, und die Zwischenpunkte je $61^{\mathrm{mm}}{,}12$ von einander entfernt liegen.
\begin{center}
\renewcommand{\arraystretch}{1.15}
\begin{tabular}{|r|r||r|r||r|r|}\hline
0 & 11,4690 & 7 & 8,0232 & 14 & 7,0562\\
1 & 9,9611 & 8 & 7,8652 & 15 & 6,9561\\
2 & 9,2551 & 9 & 7,7143 & 16 & 6,8793\\
3 & 8,8605 & 10 & 7,5695 & 17 & 6,8388\\
4 & 8,5906 & 11 & 7,4305 & 18 & 6,8507\\
5 & 8,3763 & 12 & 7,2974 & 19 & 6,8771\\
6 & 8,1917 & 13 & 7,1723 & 20 & 6,3527\\\hline
\end{tabular}
\end{center}
\clearpage\endgroup
